\chapter{Digoxin Level}

\section*{What Is This Test?}
The digoxin level (serum digoxin concentration) measures the amount of digoxin, a cardiac glycoside, in the blood. Digoxin is derived from the foxglove plant (Digitalis lanata) and is used to treat atrial fibrillation (for rate control), atrial flutter, and systolic heart failure with reduced ejection fraction (HFrEF). Digoxin exerts its therapeutic effect by inhibiting the sodium-potassium ATPase (Na+/K+-ATPase) pump on cardiac myocytes, increasing intracellular sodium, which decreases the sodium-calcium exchanger activity, thereby increasing intracellular calcium, enhancing cardiac contractility (positive inotropy). Digoxin also enhances vagal tone (parasympathetic activity), slowing conduction through the atrioventricular (AV) node (negative chronotropy). Digoxin has a narrow therapeutic index (therapeutic range 0.5--2.0 ng/mL), and toxicity (levels >2.0--3.0 ng/mL) can be life-threatening: cardiac arrhythmias (bidirectional ventricular tachycardia, atrial tachycardia with AV block, ventricular fibrillation), neurological symptoms (confusion, visual disturbances including yellow-green halos and chromatopsia), and gastrointestinal symptoms (nausea, vomiting, anorexia). Therapeutic drug monitoring of digoxin is essential because its pharmacokinetics are affected by renal function (digoxin is 70--80\% renally excreted), age, drug-drug interactions (amiodarone, verapamil, quinidine, macrolide antibiotics increase digoxin levels), and electrolyte disturbances (hypokalemia, hypomagnesemia, and hypercalcemia increase myocardial sensitivity to digoxin toxicity).

\section*{Alternative Names}
Digoxin Level; Serum Digoxin; Lanoxin Level; Digitalis Level; Digoxin Concentration; Digoxin Blood Test

\section*{Principle}
Digoxin is measured by competitive chemiluminescent immunoassay on automated platforms. Digoxin in the patient sample competes with a labeled digoxin derivative (acridinium ester or ruthenium complex) for binding to a limited quantity of anti-digoxin monoclonal or polyclonal antibodies. After incubation, separation, and washing, the bound labeled-digoxin signal is inversely proportional to the digoxin concentration. Digibind (digoxin immune Fab, the antidote for digoxin toxicity) interferes with digoxin immunoassays: after Digibind administration, both free and Fab-bound digoxin are measured, and serum digoxin levels are falsely elevated (typically >10--20 ng/mL). Free (unbound) digoxin levels can be measured by ultrafiltration or specialized methods when Digibind has been administered. Endogenous digoxin-like immunoreactive factors (DLIFs) are endogenous substances found in neonates, pregnant women, and patients with renal failure, liver disease, and volume expansion that cross-react with some digoxin immunoassays, producing falsely elevated digoxin levels in patients who are not taking the drug.

\section*{History}
Digoxin and other cardiac glycosides from Digitalis purpurea (purple foxglove) were first described by William Withering in ``An Account of the Foxglove and Some of Its Medical Uses'' (1785), documenting the use of foxglove leaf extract for the treatment of dropsy (edema, now recognized as heart failure). Digoxin was isolated and purified in 1930. The Na+/K+-ATPase pump as the molecular target of digoxin was established by Jens Christian Skou (Nobel Prize, 1997). Digoxin radioimmunoassay was developed in the 1960s, enabling therapeutic drug monitoring. Digoxin immune Fab (Digibind), a specific antidote for digoxin toxicity, was developed in the 1970s--1980s. The DIG trial (1997) demonstrated that digoxin does not reduce overall mortality in heart failure patients but does reduce hospitalizations for worsening heart failure \cite{dig1997effect}.

\section*{Why This Test Is Ordered}
Monitoring of digoxin therapy to ensure levels are within the therapeutic range (0.5--2.0 ng/mL, with the lower end [0.5--0.9 ng/mL] preferred for heart failure and 0.5--1.0 ng/mL preferred for atrial fibrillation per current guidelines) \cite{january2014aha}. Diagnosis of digoxin toxicity in a patient with clinical symptoms (nausea, vomiting, confusion, visual disturbances, cardiac arrhythmias). Assessment of digoxin level after a change in renal function, addition of an interacting drug, or suspected overdose. Monitoring after Digibind (digoxin immune Fab) administration.

\section*{Primary vs Secondary Test}
This is a primary therapeutic drug monitoring test.

\section*{How This Test Is Performed}
Blood sample in a serum separator tube (SST, gold-top) or red-top tube. The sample must be drawn at least 6--8 hours after the last oral dose (and preferably 12--24 hours after, representing a trough level), because digoxin requires 6--8 hours for tissue distribution to equilibrate; levels drawn earlier are falsely elevated (post-distributional phase). Timing of collection relative to dose is critical.

\section*{What Preparation Is Needed}
No fasting is required. Timing of the blood draw relative to the last oral dose is the single most important preparation step: the specimen must be a trough sample drawn at least 6--8 hours after the last dose and preferably 12--24 hours after it, because levels obtained earlier are falsely elevated during the post-distributional phase. The clinician should document the time of the last dose and all concurrent medications, since interacting drugs (amiodarone, verapamil, quinidine, macrolide antibiotics, itraconazole) raise digoxin levels. Steady-state levels are not interpretable until digoxin has been given at a constant dose for 5--7 days (digoxin half-life is 36--48 hours with normal renal function and longer with renal impairment) \cite{burton2006applied}.

\section*{Contraindications and Precautions}
There are no absolute contraindications to venipuncture for digoxin level measurement. Interpretive cautions: (1) endogenous digoxin-like immunoreactive factors in neonates, pregnant women, and patients with renal or hepatic failure or volume expansion can cause falsely elevated levels. (2) After Digibind (digoxin immune Fab) administration, standard immunoassays measure total (free plus Fab-bound) digoxin and produce markedly elevated levels that do not reflect toxicity; free digoxin measured by ultrafiltration is required for monitoring. (3) Hypokalemia, hypomagnesemia, and hypercalcemia increase myocardial sensitivity to digoxin, so toxicity can occur at levels within the therapeutic range. (4) Samples drawn too soon after a dose are falsely elevated and should not guide dosing. (5) In chronic therapy, toxicity typically occurs at levels above 2.0--3.0 ng/mL, whereas acute overdose produces higher levels.

\section*{Risks}
Minimal risk. Slight pain or bruising at the venipuncture site. Rarely: hematoma, infection, or vasovagal syncope.

\section*{What Happens After the Test}
Results are typically available within a few hours to 1 day. The trough level is interpreted against the therapeutic range (0.5--2.0 ng/mL) together with serum potassium, magnesium, calcium, and renal function. Subtherapeutic levels may prompt a dose increase with a repeat trough after 5--7 days; levels in the toxic range require holding the dose and correcting electrolyte disturbances, and severe toxicity or marked hyperkalemia may warrant digoxin immune Fab and cardiac monitoring.

\section*{Understanding Results}

\subsection*{Below Therapeutic Range (Subtherapeutic, <0.5 ng/mL)}
Inadequate dosing, non-adherence, malabsorption, or drug interactions increasing digoxin clearance (rifampin, certain antacids and cholestyramine reduce absorption). Minimal therapeutic effect --- inadequate heart rate control in atrial fibrillation or lack of inotropic benefit in heart failure.

\subsection*{Above Therapeutic Range (Toxic, >2.0 ng/mL)}
Digoxin toxicity. ECG findings: bidirectional ventricular tachycardia (pathognomonic but rare), atrial tachycardia with AV block, sinus bradycardia, AV block, ventricular ectopy \cite{goldberger2018clinical}. Clinical symptoms: nausea, vomiting, anorexia (GI symptoms are the earliest), confusion, delirium, visual disturbances (yellow-green chromatopsia, blurred vision, halos around lights). Hyperkalemia (serum potassium >5.0--5.5 mEq/L) in acute digoxin toxicity is a marker of severe poisoning (inhibition of Na+/K+-ATPase in skeletal muscle releases potassium) and is a predictor of mortality (serum potassium >5.5 mEq/L is an indication for Digibind therapy) \cite{bauman2006mechanisms}.

\section*{Normal Results}
Therapeutic range: 0.5--2.0 ng/mL (0.6--2.6 nmol/L). Heart failure: 0.5--0.9 ng/mL (lower end of the range is associated with reduced mortality in chronic heart failure) \cite{yancy2013accf}. Atrial fibrillation (rate control): 0.5--1.0 ng/mL. Toxicity: >2.0 ng/mL (chronic toxicity); >3.0--4.0 ng/mL (severe toxicity). Therapeutic drug monitoring must use a trough sample (drawn at least 6--8 hours after dose, ideally 12--24 hours for steady state).

\section*{Follow-up Tests}
Serum electrolytes (potassium, magnesium, calcium) --- hypokalemia, hypomagnesemia, and hypercalcemia exacerbate digoxin toxicity. Renal function (serum creatinine, BUN, eGFR) --- digoxin is renally excreted. ECG (electrocardiogram) for arrhythmias. Digoxin level after Digibind administration (note: standard immunoassays measure total [free + Fab-bound] digoxin and will be markedly elevated after Digibind). Free digoxin level by ultrafiltration if monitoring is required after Digibind.

\section*{References}
\bibliographystyle{unsrtnat}
\bibliography{references}

\clearpage
